\scanpage{176}
% L1–6
On se fixe un $\lambda\in k$, et on
travaille avec les extensions $\lambda$-scindées.
\deleted{On appelle} \unsure{se donné|se donne|Si on se donne} \unsure{de|des} telles
\diagram{\begin{tikzcd}[ampersand replacement=\&,column sep=small]
0\arrow[r,dash]\&E_i^0\arrow[r,dash,"\mathrm{inj}_i"]\&
 E_i\arrow[r,dash]\arrow[l,bend left=30,"u_i"']\& F_i\arrow[r,dash]\&0
\end{tikzcd}}
\[
u_i\circ\mathrm{inj}_i=\lambda\unsure{\xi|\epsilon|{}}\,\mathrm{id}_{E_i^0}
\qquad \lambda\ \text{\uuline{indép. de} }i
\]
et un module $B$\unsure{,} sur $k$. \deleted{Une}

% L7–12
On appelle $\lambda$-accouplement à valeurs
dans $B$, une couple
\[
\begin{aligned}
\prod E_i &\mathrel{\mathop{\text{---}}\limits^{f}} B\\
\prod E_j^0 &\mathrel{\mathop{\text{---}}\limits^{f^0}} B
\end{aligned}
\]
d'applications \deleted{\unsure{bi}} \inserted{multi}linéaires, satisfaisant
\[
f(x_1,\ldots,x_i+\alpha_i,\ldots x_n)
=f(x_1,\ldots x_i,\ldots x_n)+f^0(u_1x_1),\ldots,\alpha_i,\ldots,u_n(x_n))
\]
% L13–19
Prenant les $x_i$ dans $E_i^0$, on trouve
\[
f(\alpha_1,\ldots,\alpha_i,\ldots\alpha_n)=\lambda^{n-1}f^0(\alpha_1,\ldots,\alpha_i,\ldots\alpha_n)
\]
(on a pris $x_j=\alpha_j$ pour $j=1,\ldots\hat i,\ldots n$, et
$x_i=0$). On a une solution du
pb universel, noté $\mathop{*}\limits_{i,\lambda}E_i$ ($E_1*_\lambda E_2$ etc)
l'indice $\lambda$ rappelant que c'est relatif à un
choix de $\lambda$-trivialisations des extensions
envisagées. L'\unsure{espace|expression} $\mathop{*}\limits_{\unsure{i|j},\lambda}E_i$ se construit
% archive sheet number: 175
